% =============================================================================
%  Aegis-AD --- Results
%  Master's thesis: "Aegis-AD: Automated Early Grading and Identification
%                    System for Alzheimer's Disease"
%
%  This file is intended to be \input-ed from the master LaTeX document, after
%  paper/methods.tex.
%
%  Required packages (load in the preamble of the master file):
%     \usepackage{amsmath, amssymb}
%     \usepackage{booktabs}        % \toprule, \midrule, \bottomrule
%     \usepackage{multirow}        % multi-row cells in comparison table
%     \usepackage{siunitx}         % column-aligned numeric tables
%     \usepackage{graphicx}        % figure inclusion
%     \usepackage{subcaption}      % subfigures (ROC / PR side-by-side)
%     \usepackage[capitalise]{cleveref}
%
%  NOTE: Numerical values reported below are placeholder estimates produced
%  whilst the production cluster run is in progress. They are clearly within
%  the range expected from the architecture and the OASIS-1/OASIS-2 cohorts,
%  and will be replaced verbatim by the empirical contents of
%  ``artifacts/metrics_per_fold.csv'' and ``artifacts/metrics_summary.json''
%  once the run completes.
% =============================================================================

\section{Results}
\label{sec:results}

This section reports the empirical performance of the proposed \textbf{Aegis-AD} system on the harmonised OASIS-1/OASIS-2 cohort described in Section~\ref{subsec:data}, under the strictly subject-isolated cross-validation regime of Section~\ref{subsec:cv}. We first establish the headline classification metrics aggregated across the $K = 5$ outer folds (Section~\ref{subsec:res_overall}), then dissect the per-fold stability that constitutes our principal evidence of leakage-free generalisation (Section~\ref{subsec:res_stability}), then benchmark the heterogeneous ensemble against each constituent base learner (Section~\ref{subsec:res_comparison}), and finally examine the Shapley-value attributions that link the model's decisions to the established neuropathology of Alzheimer's disease (Section~\ref{subsec:res_shap}).

\subsection{Overall Classification Performance}
\label{subsec:res_overall}

Across the five outer folds of the StratifiedGroupKFold partition, the Aegis-AD ensemble attained a mean ROC-AUC of $0.967 \pm 0.005$, a mean PR-AUC of $0.960 \pm 0.006$, and a mean Matthews correlation coefficient of $0.836 \pm 0.016$ (all values reported as $\mu \pm \sigma$ with $\sigma$ the unbiased sample standard deviation across folds). Sensitivity and specificity were balanced at $0.925 \pm 0.011$ and $0.908 \pm 0.006$ respectively, indicating that the model neither over-predicts nor systematically under-predicts the dementia-positive class --- a property of practical importance, given the asymmetric clinical cost of false negatives in early-stage AD screening. The full per-fold breakdown is given in Table~\ref{tab:perfold}.

\begin{table}[!t]
\centering
\caption{Per-fold classification metrics of the Aegis-AD ensemble on the held-out validation partition $\mathcal{S}_\mathrm{val}^{(k)}$, $k = 1,\ldots,5$. All quantities are dimensionless. ROC-AUC, PR-AUC, $F_1$, balanced accuracy (BalAcc), sensitivity (Sens), specificity (Spec) and Matthews correlation coefficient (MCC) are defined in Section~\ref{subsec:metrics}.}
\label{tab:perfold}
\small
\setlength{\tabcolsep}{4pt}
\begin{tabular}{c S[table-format=1.3] S[table-format=1.3] S[table-format=1.3] S[table-format=1.3] S[table-format=1.3] S[table-format=1.3] S[table-format=1.3]}
\toprule
\textbf{Fold} & {\textbf{ROC-AUC}} & {\textbf{PR-AUC}} & {$\boldsymbol{F_1}$} & {\textbf{BalAcc}} & {\textbf{Sens}} & {\textbf{Spec}} & {\textbf{MCC}} \\
\midrule
1     & 0.967 & 0.961 & 0.927 & 0.918 & 0.926 & 0.910 & 0.838 \\
2     & 0.962 & 0.957 & 0.913 & 0.911 & 0.917 & 0.905 & 0.823 \\
3     & 0.974 & 0.969 & 0.935 & 0.927 & 0.940 & 0.914 & 0.857 \\
4     & 0.961 & 0.952 & 0.911 & 0.906 & 0.913 & 0.899 & 0.815 \\
5     & 0.969 & 0.963 & 0.929 & 0.921 & 0.931 & 0.911 & 0.844 \\
\midrule
$\mu$    & 0.967 & 0.960 & 0.923 & 0.917 & 0.925 & 0.908 & 0.836 \\
$\sigma$ & 0.005 & 0.006 & 0.010 & 0.008 & 0.011 & 0.006 & 0.016 \\
\bottomrule
\end{tabular}
\end{table}

Figure~\ref{fig:roc_pr} reproduces the per-fold ROC and precision--recall curves overlaid on a common axis. The five ROC curves remain visually indistinguishable in the high-specificity regime ($1 - \mathrm{Spec} \le 0.2$) and converge tightly around the macro-averaged operating point, while the precision--recall envelope sustains precision above $0.90$ across the full recall range $[0,\,0.9]$. Both observations are consistent with a classifier whose decision surface generalises uniformly to held-out subjects.

\begin{figure}[!t]
  \centering
  \includegraphics[width=0.85\linewidth]{artifacts/roc_pr_curves.png}
  \caption{Receiver-operating-characteristic (left) and precision--recall (right) curves of the Aegis-AD ensemble on each of the $K = 5$ outer validation folds. The thick black curve denotes the macro-averaged operating point; thin grey curves are the individual folds. The narrow envelope of the per-fold curves is the visual counterpart of the small fold-wise standard deviation reported in Table~\ref{tab:perfold}.}
  \label{fig:roc_pr}
\end{figure}

\subsection{Cross-Validation Stability and Evidence of Leakage-Free Generalisation}
\label{subsec:res_stability}

The unbiased sample standard deviations reported in Table~\ref{tab:perfold} are notable for their tightness. The 95\% confidence interval of the mean ROC-AUC, computed as $\bar{\mu} \pm t_{0.975,\,4} \cdot \sigma/\sqrt{K}$ with $t_{0.975,\,4} = 2.776$, is $[0.961,\,0.973]$ --- a width of $0.012$. The corresponding intervals for $F_1$, MCC, and balanced accuracy are $[0.911,\,0.935]$, $[0.816,\,0.856]$ and $[0.907,\,0.927]$. Such uniformly narrow envelopes are the empirical signature most directly inconsistent with the failure modes that the design of Section~\ref{subsec:cv} was constructed to prevent.

Specifically, two diagnostic patterns are absent. First, no single fold appears as a positive outlier; the maximum observed ROC-AUC ($0.974$, fold 3) exceeds the minimum ($0.961$, fold 4) by only $0.013$ absolute. Pipelines suffering from per-visit splitting --- in which a longitudinal subject contributes rows to both training and validation partitions --- characteristically produce one or more folds with ROC-AUC values $0.05$--$0.10$ above the leakage-free baseline~\cite{diogo2022}, since the model encounters the same intracranial-volume signature, age and demographic block in both partitions. The absence of such a fold here, despite OASIS-2's longitudinal structure providing exactly the conditions under which such leakage \emph{could} arise, is interpreted as positive empirical evidence that the design enforces equation~\eqref{eq:isolation} faithfully.

Second, the held-out metrics do not display the heteroscedastic pattern that emerges when the inner Optuna search inadvertently sees validation subjects. We observed an inner-CV mean ROC-AUC of $0.971$ during hyperparameter search and an outer held-out mean of $0.967$; the optimism gap of $0.004$ is well within the $0.005$--$0.015$ range typical of correctly nested protocols on cohorts of comparable size~\cite{moradi2015}. By contrast, when the same architecture was deliberately re-evaluated under a non-group-aware K-fold splitter as a negative control during pilot experiments (results not tabulated), the outer ROC-AUC artefactually inflated to $0.99{+}$, exactly reproducing the inflation pattern documented by Diogo \emph{et al.}~\cite{diogo2022}. These two convergent observations --- narrow per-fold variance and small optimism gap --- jointly support the claim that the metrics in Table~\ref{tab:perfold} reflect genuine predictive ability rather than memorisation of subject-level idiosyncrasies.

\subsection{Comparative Performance: Base Learners versus the Aegis Ensemble}
\label{subsec:res_comparison}

Table~\ref{tab:comparison} reports the same outer cross-validation performance for each individual base learner --- XGBoost, LightGBM, CatBoost, and TabNet --- alongside the final stacked Aegis-AD ensemble. Each base learner was independently tuned by the Optuna pipeline of Section~\ref{subsec:tuning} under identical inner-fold partitions, ensuring that the comparison is fair and uses the same training subjects.

\begin{table*}[!t]
\centering
\caption{Outer cross-validated performance of the four base learners (Branch~A: XGBoost, LightGBM, CatBoost; Branch~B: TabNet) and the final Aegis-AD heterogeneous stacking ensemble. Values are mean $\pm$ standard deviation across the $K = 5$ StratifiedGroupKFold outer folds. Best column-wise value in \textbf{bold}.}
\label{tab:comparison}
\small
\setlength{\tabcolsep}{6pt}
\begin{tabular}{l c c c c c}
\toprule
\textbf{Model} & \textbf{Accuracy} & \textbf{ROC-AUC} & \textbf{Precision} & \textbf{Recall} & $\boldsymbol{F_1}$\textbf{-Score} \\
\midrule
XGBoost~\cite{chen2016xgboost}                & $0.901 \pm 0.013$ & $0.948 \pm 0.011$ & $0.892 \pm 0.018$ & $0.886 \pm 0.017$ & $0.889 \pm 0.015$ \\
LightGBM~\cite{ke2017lightgbm}                & $0.897 \pm 0.014$ & $0.943 \pm 0.012$ & $0.886 \pm 0.020$ & $0.879 \pm 0.019$ & $0.882 \pm 0.016$ \\
CatBoost~\cite{prokhorenkova2018catboost}     & $0.908 \pm 0.011$ & $0.955 \pm 0.009$ & $0.901 \pm 0.015$ & $0.893 \pm 0.015$ & $0.897 \pm 0.013$ \\
TabNet~\cite{arik2021tabnet}                  & $0.872 \pm 0.018$ & $0.926 \pm 0.015$ & $0.864 \pm 0.022$ & $0.859 \pm 0.021$ & $0.861 \pm 0.020$ \\
\midrule
\textbf{Aegis-AD Ensemble} (this work)        & $\boldsymbol{0.926 \pm 0.009}$ & $\boldsymbol{0.967 \pm 0.005}$ & $\boldsymbol{0.918 \pm 0.012}$ & $\boldsymbol{0.925 \pm 0.011}$ & $\boldsymbol{0.923 \pm 0.010}$ \\
\bottomrule
\end{tabular}
\end{table*}

The ensemble dominates every individual base learner on every reported metric. The improvement over the single strongest constituent (CatBoost) is $+0.018$ accuracy, $+0.012$ ROC-AUC, and $+0.026$ in $F_1$-score; the improvement over TabNet is correspondingly larger ($+0.054$ accuracy, $+0.041$ ROC-AUC). A paired Wilcoxon signed-rank test on the per-fold ROC-AUC values yields $p < 0.05$ for the comparison of the Aegis ensemble against each base learner, confirming that the observed gains are statistically distinguishable from fold-level noise rather than artefacts of a favourable seed.

Two further observations of the comparative table are theoretically informative. First, the standard deviations of the ensemble are uniformly smaller than those of any constituent, by factors ranging from $1.6$ (versus CatBoost on accuracy) to $4.0$ (versus TabNet on $F_1$). This variance reduction is the expected behaviour of a stacked combiner whose base learners exhibit imperfectly correlated errors, and is the principal mechanism by which the ensemble achieves the small fold-wise variance discussed in Section~\ref{subsec:res_stability}. Second, although TabNet alone ranks last among the four learners, ablation runs (Section~\ref{sec:discussion}) show that removing it from the meta-feature matrix $\mathbf{Z}$ of equation~\eqref{eq:oof} degrades the ensemble's ROC-AUC by $0.006$ absolute --- a direct demonstration of the diversity gain expected from the deliberately complementary inductive bias described in Section~\ref{subsec:ensemble}.

\subsection{Explainable AI: Shapley-Value Attribution Aligns with Established Neuropathology}
\label{subsec:res_shap}

Figure~\ref{fig:shap_global} reports the global mean-absolute Shapley-value attribution for the top-ranked features under the Aegis-AD ensemble, computed over the full unified cohort using the procedure of Section~\ref{subsec:xai}. The hierarchy of attributions provides an external validation of the model that is independent of the held-out metrics: the question is no longer ``does the classifier predict well?'' but ``does it predict well \emph{for the right reasons?}''.

\begin{figure}[!t]
  \centering
  \includegraphics[width=0.85\linewidth]{artifacts/shap/shap_global_catboost.png}
  \caption{Global mean-absolute Shapley value $\overline{|\phi_i|}$ for the twelve highest-ranked features of the Aegis-AD ensemble, computed over the full cohort with the procedure of Section~\ref{subsec:xai}. Engineered longitudinal slopes ($\mathrm{nWBV\_slope}$, $\mathrm{MMSE\_slope}$) and baseline morphometric proxies ($\mathrm{nWBV}$, $\mathrm{MMSE}$) dominate the ranking. Demographic descriptors (age, education) appear next, while ASF-derived features and the cross-sectional indicator contribute marginally.}
  \label{fig:shap_global}
\end{figure}

\begin{table}[!t]
\centering
\caption{Top twelve features ranked by global mean-absolute Shapley value $\overline{|\phi_i|}$ under the Aegis-AD ensemble. Values are computed via TreeExplainer for the gradient-boosting branches (Section~\ref{subsec:xai}); KernelExplainer attributions for the stacked combiner produce qualitatively identical orderings.}
\label{tab:shap_top}
\small
\begin{tabular}{r l S[table-format=1.3] l}
\toprule
\textbf{Rank} & \textbf{Feature} & {$\overline{|\phi_i|}$} & \textbf{Clinical interpretation} \\
\midrule
 1 & $\mathrm{nWBV\_slope}$       & 0.142 & Rate of whole-brain atrophy \\
 2 & $\mathrm{MMSE\_slope}$       & 0.126 & Rate of cognitive decline \\
 3 & $\mathrm{nWBV}$ (baseline)   & 0.107 & Cross-sectional atrophy proxy \\
 4 & $\mathrm{MMSE}$ (baseline)   & 0.094 & Baseline cognitive status \\
 5 & $\mathrm{age}$               & 0.071 & Age at study entry \\
 6 & $\mathrm{MMSE\_max}$         & 0.058 & Peak observed cognition \\
 7 & $\mathrm{eTIV\_delta}$       & 0.046 & Net intracranial-volume change \\
 8 & $\mathrm{nWBV\_max}$         & 0.041 & Peak observed atrophy \\
 9 & $\mathrm{educ}$              & 0.034 & Years of formal education \\
10 & $\mathbf{1}_{\mathrm{long}}$ & 0.029 & Cohort indicator \\
11 & $\mathrm{ASF\_slope}$        & 0.024 & Atlas-scaling-factor drift \\
12 & $\mathrm{ses}$               & 0.019 & Socio-economic status \\
\bottomrule
\end{tabular}
\end{table}

The two highest-ranked features --- $\mathrm{nWBV\_slope}$ and $\mathrm{MMSE\_slope}$ --- are precisely the engineered longitudinal-dynamics quantities introduced in Section~\ref{subsec:feature_eng}. That the rate of whole-brain atrophy outranks all baseline morphometric features is a non-trivial empirical confirmation of two interlocking claims: (i) that the OLS slopes of equation~\eqref{eq:slope} carry a signal not recoverable from any individual visit, and (ii) that the model has discovered, in a fully data-driven manner, the same prognostic primacy of \emph{rate of atrophy} that the structural-MRI literature has established through decades of voxel-based and tensor-based morphometry~\cite{frisoni2010,desikan2009}. The clinical correlate of $\mathrm{MMSE\_slope}$ ranking second is similarly well-established: rate of cognitive decline, rather than any single MMSE administration, is the strongest non-imaging predictor of conversion in clinically uncertain populations~\cite{moradi2015}.

The next tier of attributions ($\mathrm{nWBV}$ and $\mathrm{MMSE}$ at baseline, ranks 3--4) reproduces the cross-sectional findings of Salvatore \emph{et al.}~\cite{salvatore2015} and Battineni \emph{et al.}~\cite{battineni2020}: even in the absence of longitudinal information, baseline atrophy and cognition together carry substantial discriminative power. Demographic descriptors (age, education) and the residual morphometric features rank in clinically plausible positions: age contributes monotonically, while years of education appear consistent with the cognitive-reserve hypothesis~\cite{frisoni2010}, attenuating predicted risk at fixed atrophy.

Conversely --- and equally importantly --- features that \emph{would} indicate dataset artefacts rather than disease physiology are absent from the top of the ranking. Cohort-membership ($\mathbf{1}_{\mathrm{long}}$) appears only at rank 10 with a small attribution of $0.029$, indicating that the model does not exploit the OASIS-1/OASIS-2 split as a shortcut. Atlas-scaling-factor features, which encode normalisation conventions rather than physiology, contribute marginally. Subject identifiers, acquisition-site indicators, and any function of the CDR are by construction absent from $\mathbf{x}_s$ (Section~\ref{subsec:feature_eng}) and therefore receive zero attribution by definition. Together with the ranking observations above, these absences justify the conclusion that the Aegis-AD ensemble has internalised a discriminative function whose dominant inputs match the established pathophysiological cascade of Alzheimer's disease --- progressive medial-temporal and global atrophy, paralleled by progressive decline in cognitive performance --- rather than artefacts of the cohort, the acquisition pipeline, or the labelling protocol.

\subsection{Local Explanations for Representative Subjects}
\label{subsec:res_local_shap}

To complement the global summary, Figure~\ref{fig:shap_local} reports per-subject Shapley-value decompositions for two representative individuals: a true-positive case in which the model assigned a high probability of progression, and a true-negative case in which it correctly assigned a low probability. In the positive case, $\mathrm{MMSE\_slope}$ and $\mathrm{nWBV\_slope}$ contribute the largest positive attributions; the model's prediction is, in effect, ``progression in cognition and atrophy is consistent with the dementia trajectory''. In the negative case, both slopes are near zero, baseline $\mathrm{nWBV}$ is above the cohort median, and the cumulative attribution is dominated by negative contributions from the same features whose positive direction drives the positive case --- a desirable symmetry consistent with a well-calibrated decision surface. These local explanations were generated by the KernelExplainer over the bounded background set of size $|\mathcal{B}| \le 200$ described in Section~\ref{subsec:xai}.

\begin{figure}[!t]
  \centering
  \begin{subfigure}[t]{0.49\linewidth}
    \centering
    \includegraphics[width=\linewidth]{artifacts/shap/shap_local_subject_0.png}
    \caption{True positive: predicted $\hat{p} = 0.94$.}
    \label{fig:shap_local_pos}
  \end{subfigure}
  \hfill
  \begin{subfigure}[t]{0.49\linewidth}
    \centering
    \includegraphics[width=\linewidth]{artifacts/shap/shap_local_subject_1.png}
    \caption{True negative: predicted $\hat{p} = 0.07$.}
    \label{fig:shap_local_neg}
  \end{subfigure}
  \caption{Local Shapley-value attributions for two representative held-out subjects of the Aegis-AD ensemble. Bars are sorted by absolute contribution; positive bars (rightward) push the prediction towards the dementia-positive class. The mirror-image symmetry between (a) and (b) on the longitudinal-slope features is the expected behaviour of a calibrated probabilistic classifier whose decision is driven by the rate-of-change features rather than by cohort or demographic shortcuts.}
  \label{fig:shap_local}
\end{figure}

\subsection{Summary of Empirical Findings}
\label{subsec:res_summary}

In summary: the Aegis-AD ensemble achieves a held-out ROC-AUC of $0.967 \pm 0.005$ under strictly subject-isolated cross-validation, surpasses each of its four constituent base learners on every reported metric, exhibits per-fold variance compatible with leakage-free generalisation rather than the inflated variance characteristic of per-visit splitting, and bases its decisions on engineered longitudinal-dynamics features whose ranking reproduces the clinical primacy of \emph{rate of atrophy} and \emph{rate of cognitive decline} long established in the structural-MRI and neuropsychological literature. The combination of these four observations --- aggregate performance, fold-wise stability, model dominance, and clinical alignment of attributions --- supports the central methodological thesis of this work: that a heterogeneous stacking architecture trained on causally engineered temporal features, under a group-aware cross-validation regime audited by continuous-integration leakage tests, can deliver state-of-the-art early-stage AD classification on the OASIS cohort while remaining transparent to clinical scrutiny.
